% AI-written, AI cross-reviewed; no human sign-off; not a claim that RH is proved
% RH Notes blueprint source excerpts for Z2.K7
% Verbatim excerpts in the notation of the parent blueprint. The labelled editions and source ranges appear below.
% See the unit YAML for the current proof status. Earlier-round passages retain their original wording.
\section*{Z2.K7. Full integral homology torsion}

% Parent source SHA-256: 1af64922b98ed2310744411b1f18b25cc426c97220646321471836e1a3a409d6
% blueprint/ch_P.tex lines 652 through 727; prop:P-k-num
\begin{proposition}[Absorbing the two torsion orders]\label{prop:P-k-num}
\lean{UpperTheory.V.ZK.sNum}
\uses{prop:P-k-layer,prop:P-k-count}
There is a single $C_\Sigma$ such that $\Sigma(L)=s_0<\infty$ and
$s_0\log\ell\le K_{C_\Sigma}(E,n,D)$ in Proposition~\ref{prop:P-k-main}.
\end{proposition}
\begin{proof}
Section~\ref{sec:P-r3-K247} proves the class-number estimate and the numerical step of K7. Its geometric and homological bounds are stated there with their proof sources. For every totally complex $M$, every $E\ge1$
and every $i\le2E+1$ the integral torsion in rank $4E+6$ is finite and
$\log T_i\le K_{C_0}(E,n',D')$ with one absolute $C_0$. The project
obtains $C_0=24$ by combining a torsion estimate for the arithmetic group,
a cell-orbit bound and the numerical comparison \nolinkurl{exponent_le}
(\nolinkurl{torsAll}, \nolinkurl{sTors}, \path{ExponentLe.lean:6,70--125}).
Here \nolinkurl{sTors} passes only the logarithmic inequality;
finiteness is the other conjunct of \nolinkurl{torsAll}, supplied by
\nolinkurl{KT.tors_le} (\path{HermiteNormalForms/Assembly.lean:25--30}).
The auxiliary class-number bound used below is
$\lfloor\log_2 h_M\rfloor\le4n'\log(2D')$. In \nolinkurl{gens_le}
(\path{NumThree.lean:186--208}) this follows by setting the prime equal
to $2$ in \nolinkurl{numGens} (\path{GeneratorCountTuples/NumGens.lean:6--45}),
which invokes \nolinkurl{log_classNumber_le} (\path{IdealCounting/ClassC.lean:63}).
The factors $T_i$ are taken in degrees $2E$ and $2E+1$, so both meet the
homology cutoff, including its top endpoint.

By the preceding propositions, $\Sigma(L)\le s$ in the extended
naturals, hence $\Sigma(L)$ is finite before one takes its integer value
$s_0$. Also $1\le\ell^s\le T_{2E+1}T_{2E}\ell^X$. In particular both
torsion cardinalities are strictly positive independently of the convention
for the cardinality of an infinite type. We may take logarithms:
\[
  s_0\log\ell\le s\log\ell
  \le \log T_{2E+1}+\log T_{2E}+X\log\ell .
\]
See \nolinkurl{sNum} (\path{NumThree.lean:241--284}). Replacing extended
length by its natural part prematurely would also assign the natural
value zero to infinity.

Put $u=2En'\ge2$, $Y=\log(2D')\ge1/2$. Since $\ell,c\le\ell^c\le2n'$,
both are at most $u$, as are $v$ and $\log\ell$. From $2E+1<r_2(M)$
and $2n'\le u$ one gets
\[
 b\le u+2u^2,\qquad g\le u+2uY+2 .
\]
For the first inequality bound the product in $b$ by $n'(3n')$.
For the second insert the class-number estimate in the definition of
$g$. Hence $b+v\le3u^2$, $g\le6uY$, and $c+v\le4uY\le u^3Y$. Thus
\[
 X\log\ell\le19u^4Y\le u^9Y
\]
since $u^5\ge32$. For $A=u^{\max(C_0,9)}Y\ge1$ we conclude
$s\log\ell\le2e^A+A\le e^{4A}$. Indeed $A\le e^A$ and $3\le e^{3A}$.
As $4\le u^2$, this gives $K_{C_1}(E,n',D')$ with
$C_1=\max(C_0,9)+2$. The two comparisons are \nolinkurl{num_poly} and
\nolinkurl{num_three} (\path{NumThree.lean:6--76}).

It remains to transfer the real bound from $M$ down to $L$.
Write $x=2En\ge2$ and $Y_0=\log(2D)\ge1/2$.
The degree bound gives $d\le16En=8x$, $n'=dn\le4x^2$, and
$2En'=dx\le8x^2\le x^5$. Put $w=\ell^c\le2n'\le8x^2$. Since $D,D',w$
are positive, the discriminant inequality gives
\[
 \log(2D')\le \log2+d\log D+n'\log w
             \le Y_0+8xY_0+32x^4
             \le73x^4Y_0\le x^{11}Y_0 .
\]
Here $\log w\le w$ and $x^7\ge128$. These estimates also apply at $D=1$
because of the factor $2$ inside the logarithms. Raising $2En'\le x^5$
to the nonnegative exponent $C_1$ and multiplying proves
\[
  K_{C_1}(E,n',D')\le K_{5C_1+11}(E,n,D).
\]
This proves the proposition with $C_\Sigma=5C_1+11$
(\nolinkurl{num_down}, \nolinkurl{sNum}, \path{NumThree.lean:79--184,211--219,364--378}).
\emph{Status: complete modulo Propositions~\ref{prop:P-k-layer}
and~\ref{prop:P-k-count}.}
\end{proof}

% blueprint/ch_P.tex lines 1696 through 1951; prop:P-r4-k7
\begin{proposition}[An all-prime bound for the face complex]
\label{prop:P-r4-k7}
\lean{UpperTheory.V.KT.tors_le}\quad\lean{UpperTheory.V.KT.orbnum}
Let $M$ be totally complex, $n=[M:\mathbf Q]$, $D=|\operatorname{disc}M|$,
$N\ge1$, $d=nN$, and $G=\mathrm{GL}_N(\mathcal O_M)$.
For every $m\ge0$ the entire integral torsion group in $H_m(G,\mathbf Z)$
is finite of order at most
\[
 (m+2)^A s^{m A s^m},\qquad
 s=3^{dN},\quad A=(2Dd)^{12(m+1)d^3N}.
\]
Together with the two numerical estimates of Section~\ref{sec:P-r3-K247}
this closes K7 at exponent $24$ (\nolinkurl{torsAll},
\path{ClassCounting/ExponentLe.lean:70--117}).
\end{proposition}
\begin{proof}
\emph{The complex and exactness.}
Take one complex embedding from each conjugate pair. Let $V$ be the
real vector space of tuples $h=(h_v)_v$ of Hermitian $N$-by-$N$
matrices, $T$ its cone of positive semidefinite tuples, and
$\Lambda=\mathcal O_M^N$. Write $y_v$ for the embedded vector and put
\[
 q_y(h)=\sum_v y_v^*h_v y_v,\quad
 P=\{h:q_y(h)\ge1\ (0\ne y\in\Lambda)\},\quad
 U=\{x\in V^*:x(k)>0\ (0\ne k\in T)\}.
\]
The lattice embeds as a full real lattice in $\prod_v\mathbf C^N$
of dimension $d$. In an integral basis the full embedding matrix
(including conjugates) has determinant of absolute value $\sqrt D$,
since its transposed product is the trace matrix. Taking real and
imaginary parts contributes $2^{-n/2}$ per coordinate. Thus its
covolume is $(2^{-n/2}\sqrt D)^N$.
Scaling the identity form puts it in $P$ since a full lattice has a
positive shortest length. The trace functional belongs to $U$
(\nolinkurl{gramq}, \nolinkurl{gramdet}, \nolinkurl{psdcone},
\path{HermiteNormalForms/RealVec.lean:6},
\path{HermiteNormalForms/Gramdet.lean:33--49},
\path{HermiteNormalForms/Psdcone.lean:6}).

Every $h\in P$ is positive definite. A negative real direction would
give a negative rational direction in lattice coordinates and then a
negative integral one. For a semidefinite form with nonzero kernel
choose a nonzero real vector $v$ in that kernel. In integral
coordinates the multiples of $v$ modulo $\mathbf Z^d$ have a positive
integer difference $kv$ arbitrarily close to an integer vector, by
subdividing the unit cube into boxes. To ensure the integer vector is
nonzero, first replace $v$ by a scalar multiple of norm $>2$.
The quadratic value of the close integer vector is the quadratic value
of its error, contradicting the lower bound. At each $h\in P$ there
is consequently a common positive lower eigenvalue bound in a
neighborhood of $h$. Only finitely many $y$ can have $q_y\le1$ anywhere
in that neighborhood (\nolinkurl{grampd}, \nolinkurl{rypos},
\nolinkurl{locfin}, \path{HermiteNormalForms/Grampd.lean:8},
\path{HermiteNormalForms/Rypos.lean:9}, \path{HermiteNormalForms/Locfin.lean:8}).

Call $v$ a vertex if the linear forms $q_y$ with $q_y(v)=1$ span $V^*$.
Let $C_v$ be the closed convex set of $x$ minimized at $v$ on $P$,
that is, $x(v)\le x(h)$ for all $h\in P$. The face sets $X_p$ are the
ordered $(p+1)$-tuples of these vertices with $U\cap\bigcap_j C_{v_j}$
nonempty. Repetitions are allowed, and deletion gives the face maps
(\nolinkurl{IsVert}, \nolinkurl{normalCone}, \nolinkurl{simp},
\path{ConvexNerveComplexes/NerveFaceClosed.lean:33--55};
\nolinkurl{Sx}, \path{HermiteNormalForms/SmulNeZeroLat.lean:76--91}).

These cones cover $U$ and only finitely many meet any compact
$K\subset U$. Indeed on $K\times(T\cap\{\|k\|=1\})$ evaluation
has a uniform positive minimum (empty $K$ is immediate). For a fixed
$h_0\in P$, all minimizers of all $x\in K$ lie in the closed bounded
set $P\cap\{\|h\|\le \max_{x\in K}x(h_0)/\min x(k)\}$.
Each minimum is attained. Its compact convex minimizing face has an
extreme point, obtained, for example, by successively minimizing a
basis of linear coordinates. At such a point the active $q_y$ span:
otherwise a common annihilating direction permits small perturbations
of either sign, by local finiteness, and both must still minimize.
This contradicts extremality. On the compact set of potential
minimizers, finitely many of their neighborhoods above contain
all possible active labels. Each vertex is specified by a basis drawn
from those labels, so only finitely many vertices occur
(\nolinkurl{exists_vert}, \nolinkurl{cone_finite}, \nolinkurl{vert_finite},
\path{ConvexNerveComplexes/ExistsVert.lean:5--44},
\path{ConvexNerveComplexes/ConeFinite.lean:7--37},
\path{ConvexNerveComplexes/VertFinite.lean:5}).

Here is a hand proof of the exactness used in \nolinkurl{acyc_simp}
(\path{ConvexNerveComplexes/AcycSimp.lean:9--19}); the source proves
the finite convex nerve step by hyperplane cuts
(\nolinkurl{nerve_compact}, \path{ConvexNerveComplexes/NerveCompact.lean:5--83}).
Choose a witness in $U$ for every tuple in a given finite cycle.
Their convex hull $K$ is compact, convex and contained in $U$.
The nonempty $F_v=K\cap C_v$ form a finite closed convex cover of $K$.
Its nerve equals that of the open $\delta$-neighborhoods of the $F_v$
for sufficiently small $\delta>0$: for each subfamily with empty
intersection a sequence of points at distance tending to zero from
all its members would have a common limit, by compactness. There are
only finitely many subfamilies. The thickened cover has convex
intersections and its union is the $\delta$-neighborhood of $K$, hence
convex. To see exactness with ordered tuples, take the double complex of
singular chains on intersections, summing over ordered index tuples,
with alternating deletion horizontally. Vertically each nonempty
intersection contracts to a point, leaving the ordered nerve chains.
Horizontally, for a singular simplex lying in some cover member, prepend
any chosen index whose member contains it: this contracts its augmented
index complex. The resulting chains of small simplices compute the
homology of the union, since barycentric subdivision on each finite
chain makes all simplices subordinate, and subdivision is chain
homotopic to the identity (the simplex prism supplies the homotopy).
Both computations of the double complex therefore give reduced
homology zero. The given cycle bounds already in this finite nerve,
hence in $X_\bullet$. The augmentation is onto because $U$ is nonempty.
This is the exact complex of \nolinkurl{sx_exact} and \nolinkurl{tupExact}
(\path{HermiteNormalForms/Assembly.lean:14--17},
\path{TorsionBoundsChains/TupExact.lean:58}).

\emph{Counting orbits.}
Take a tuple and a common minimizing functional $x\in U$.
At its barycenter $h$, the vectors with $q_y(h)=1$ are exactly those
minimal at every vertex. They span $M^N$ over $M$. Otherwise take a
nonzero $M$-linear functional vanishing on their span. The sum over
embeddings of its squared absolute values is a nonzero $k\in T$
annihilating all active constraints. By local finiteness $h-\epsilon k$
still lies in $P$ for small $\epsilon>0$, whereas $x(h-\epsilon k)<x(h)$.
Choose $N$ independent common minimal columns and form their integral
matrix $B$ (\nolinkurl{simp_span}, \nolinkurl{spanpsd}, \nolinkurl{sx_basis},
\path{ConvexNerveComplexes/SimpSpan.lean:5--48},
\path{HermiteNormalForms/Spanpsd.lean:6}, \path{HermiteNormalForms/ColsMinSmul.lean:13--22}).

Set $I_0=D^N d^d$, $R=N^N I_0^2$. For $h'\in P$ the real Gram matrix
$H$ of $q(h')$ in integral coordinates satisfies
\[
 \det H=2^{-d}D^N\Big(\prod_v\det h'_v\Big)^2,\qquad
 \det H\ge d^{-d}.
\]
The first identity follows from the embedding determinant above and the
real determinant $(\det h'_v)^2$ of a complex Hermitian form. For the
second, transform the lattice by its positive square root so the
minimum distance is at least one. The open cube of side $1/\sqrt d$
embeds modulo that lattice since its diameter is at most one and the
endpoints realizing one are excluded. Its volume is at most the
covolume $\sqrt{\det H}$. In particular $\prod_v\det h'_v\ge I_0^{-1}$.
This packing argument gives sufficient slack for \nolinkurl{detlow}
(\path{HermiteNormalForms/Gramlow.lean:43}, \path{HermiteNormalForms/Detlow.lean:36}).

Write $I=|N_{M/\mathbf Q}\det B|\ge1$. For any vertex in this tuple,
$J_v=B_v^* h'_v B_v$ has diagonal entries between zero and one, since
their sums over $v$ are the lengths of the chosen columns.
Hadamard gives $\det J_v\le1$; the product is
$I\prod_v\det h'_v\ge I/I_0$. Therefore $I\le I_0$ and every
$\det J_v\ge 1/I_0$. As $\operatorname{tr}J_v\le N$, the least
eigenvalue is at least $1/(I_0N^{N-1})$. Put $B'=IB^{-1}$, integral
by the adjugate identity and the fact that $I/\det B$ is integral
(the characteristic polynomial of $\det B$ gives this quotient).
If $q_y(h')=1$, then
$(B'_vy_v)^*J_v(B'_vy_v)\le I^2$ at every place. Thus all conjugates
of all coordinates of $B'y$ have modulus at most
$I\sqrt{I_0 N^{N-1}}\le R$
(\nolinkurl{nrmle}, \nolinkurl{scinv}, \nolinkurl{nearcoord},
\path{HermiteNormalForms/DetPull.lean:35},
\path{HermiteNormalForms/Scinv.lean:6}, \path{HermiteNormalForms/Nearcoord.lean:10}).

For fixed $I$, the residue matrix of $B'$ modulo $I\mathcal O_M$
determines $B$ up to left multiplication by $G$. Indeed, if $B'_1$
and $B'_2$ are congruent then $B_2 B_1^{-1}=1+B_2(B'_1-B'_2)/I$
and its inverse are integral. This leaves at most
$\sum_{I\le I_0,\, I\ge1} I^{nN^2}\le I_0^{1+dN}$ choices.
For fixed $B$, an algebraic integer whose conjugates all have modulus
at most $R$ is distinguished from all others by its class modulo
$2R+1$: the quotient of a difference by $2R+1$, if integral, has all
conjugates of modulus $<1$ and therefore has zero norm. Since $B'$ is
injective, at most $Z=(2R+1)^d$ vectors can occur as minima at a
vertex with these columns. Each vertex is determined by at most
$\dim_{\mathbf R} V\le dN$ active equations $q_y(h')=1$.
Padding an ordered choice of such equations by a blank gives at most
$(Z+1)^{dN}$ vertices for those columns. Hence the number of tuple
orbits is at most $I_0^{1+dN}(Z+1)^{dN(p+1)}$
(\nolinkurl{hnfeq}, \nolinkurl{hnf}, \nolinkurl{cntconj}, \nolinkurl{vxdet},
\nolinkurl{orb}, \path{HermiteNormalForms/Hnfeq.lean:6},
\path{HermiteNormalForms/Hnf.lean:7}, \path{HermiteNormalForms/Cntconj.lean:7},
\path{HermiteNormalForms/FiniteSubfamilySpansDual.lean:37},
\path{HermiteNormalForms/ColsMinSmul.lean:34--75}).
Put $B_0=2Dd\ge2$. Then $I_0\le B_0^{2d}$, $N^N\le B_0^d$,
$2R+1\le B_0^{5d+2}$ and $Z+1\le B_0^{(5d+2)d+1}$.
The total exponent is at most $12(p+1)d^3N$, since the first
term $2d(1+dN)$ is at most $4(p+1)d^3N$ and the remaining term
at most $8(p+1)d^3N$. This proves the asserted common bound $A$
through $p=m$ (\nolinkurl{orbnum}, \path{HermiteNormalForms/Orbnum.lean:6--88}).

\emph{Bounding stabilizers.}
The action of $G$ on forms is inverse pullback, permuting the $q_y$
and preserving the normal-cone intersection condition
(\nolinkurl{smul_mem_simp}, \path{HermiteNormalForms/SmulNeZeroLat.lean:62--75}).
The stabilizer of a tuple fixes its first positive definite vertex.
Every column of a matrix fixing this form has fixed quadratic length,
so each column ranges over a bounded subset of the integral lattice.
Thus this stabilizer is finite. On any finite subgroup, entrywise
reduction modulo $3\mathcal O_M$ is injective. Indeed if its kernel
were nontrivial it would contain an element of prime order $p$.
Write it as $1+3^aY$ with $a\ge1$ maximal. For $p\ne3$, expansion
modulo $3^{a+1}$ gives $pY=0$ modulo $3$, impossible. For $p=3$,
the expansion modulo $3^{a+2}$ gives $Y=0$ modulo $3$, again
impossible (the terms after $3^{a+1}Y$ are divisible by $3^{a+2}$).
The finite matrix ring has $3^{nN^2}=s$ elements
(\nolinkurl{stabfin}, \nolinkurl{torsfree}, \nolinkurl{mink3},
\nolinkurl{stab}, \path{HermiteNormalForms/Stabfin.lean:10},
\path{HermiteNormalForms/Torsfree.lean:35--48},
\path{HermiteNormalForms/Mink3.lean:8}, \path{HermiteNormalForms/ImportsAlt.lean:8--26}).

\emph{The integral determinant and stabilizer contributions.}
Write $C_i=\mathbf Z[X_i]$, $M_0=\mathbf Z$ and
$M_{i+1}=\ker(C_i\twoheadrightarrow M_i)$, using the exact
augmented complex above. On a transitive orbit with stabilizer $H$,
the coefficient module is $\mathbf Z[G]\otimes_{\mathbf Z[H]}\mathbf Z$.
Restricting a free right $\mathbf Z[G]$-resolution of $\mathbf Z$
to $H$ is again free and exact. Tensoring with the displayed induced
module thus computes $H_j(H,\mathbf Z)$. For $j>0$, multiplication
by $|H|$ kills this group: in homogeneous bar chains prepend every
$h\in H$ and sum. This operation is equivariant and its boundary
identity on positive degrees gives a homotopy of multiplication by
$|H|$ to zero on coinvariants. The inhomogeneous degree-$j$ chains
are free of rank $|H|^j$, so their cycles have rank at most $|H|^j$.
Their homology quotient, killed by $|H|$, has order at most
$|H|^{|H|^j}$. Summing over the at most $A$ orbits in $X_i$ gives
$|H_j(G,C_i)|\le s^{A s^j}$ for $i\le m$ and $j>0$
(\nolinkurl{shapiro}, \nolinkurl{finGroup}, \nolinkurl{permCard},
\path{TorsionBoundsChains/Shapiro.lean:6},
\path{TorsionBoundsChains/FinGroup.lean:7--19},
\path{TorsionBoundsChains/PermCard.lean:12}).

The long exact sequence for each $0\to M_{i+1}\to C_i\to M_i\to0$
maps $H_j(G,M_i)$ into $H_{j-1}(G,M_{i+1})$ with kernel an image
of $H_j(G,C_i)$. Its torsion subgroup therefore has order at most
$s^{A s^j}$ times the torsion order in the target whenever that
target's torsion is finite. After $m$ steps the last group is the
coinvariant quotient
\[
 H_0(G,M_m)=\operatorname{coker}\big((C_{m+1})_G\longrightarrow(C_m)_G\big).
\]
Here $(C_m)_G$ is free of rank at most $A$ and every column has
$\ell^1$ norm at most $m+2$, including when several faces coincide.
Choose $r$ rationally independent image columns spanning the
rational image, and $r$ coordinate rows giving a nonzero minor.
Projection to these rows is injective on that span, maps its saturated
integer lattice into $\mathbf Z^r$, and maps the chosen columns to
a lattice of index the absolute determinant of the minor. Therefore
the index of the entire image in its saturation, which equals the
cokernel torsion order, is at most that determinant. Hadamard bounds
it by $(m+2)^r\le(m+2)^A$. If $r=0$ the index is one.
This argument also works when $X_{m+1}$ has infinitely many orbits.
Multiplying all the finite bounds and using
$\sum_{j=1}^m s^j\le m s^m$ gives the proposition, including $m=0$
(\nolinkurl{cokerTors}, \nolinkurl{det_natAbs_le_prod_l1_norms},
\nolinkurl{shift}, \nolinkurl{mainFace},
\path{TorsionBoundsChains/CokerTors.lean:10},
\path{TorsionBoundsChains/DetNatAbsLeProdL1Norms.lean:5},
\path{TorsionBoundsChains/LogLeExp.lean:14--67}).
\emph{Status: K7 complete.}
\end{proof}

